% =============================================================================
% HELD BACK from the thesis on 6 Oct 2026 (to be rewritten by the author).
% Not part of the Overleaf build. Was chapter/4_conclusion.tex in full.
% =============================================================================

\chapter{Conclusions and Future Work}
\label{ch:conclusions}

\section{Conclusions}
\label{sec:conclusions}

The four research questions of Section~\ref{sec:aim} can now be answered.

\paragraph{RQ1: Can a soundtrack's emotional profile predict its film's genre?}
Yes, well above chance, though far from perfectly. On the five-genre Eerola subset, the
emotions rated by listeners reach a macro-$F_1$ of 0.428 and the emotions predicted from
audio 0.417, against 0.168 for the most-frequent baseline and 0.308 for a random guess at
the genres' base rates. On all eight genres the pipeline reaches 0.318 against a
most-frequent baseline of 0.102. Predicted emotions are as useful as the ratings they
imitate ($-0.011$, $p=0.325$), so the trained pipeline needs no manual annotation. The
caveat is the absolute level: genre labels describe films rather than clips, rare genres
have a handful of films, and the learning curves show that more films would still help.

\paragraph{RQ2a: Which emotions relate to which genres?}
Every genre that can be predicted has a readable emotional signature. Horror and Action are
defined by threat -- fear, tension and low valence -- with Horror set apart by fear; Comedy
and Drama by high valence and the absence of threat; and Adventure, the one genre without a
distinctive signature, is also the one \emph{frequent} genre the model cannot predict
(Biography and Documentary fail for lack of films). These signatures are not
an artefact of one dataset: recomputed on the 110 Blockbuster films from emotions that the
Eerola-trained model predicts for each cue, they correlate with the Eerola signatures at
$r=0.836$ across 48 genre--emotion cells, with 92\,\% agreement in sign. The caveat is that
the Blockbuster emotions are predicted, not rated, so the replication validates the
signatures indirectly.

\paragraph{RQ2b: Does the emotion intermediate beat direct audio?}
In-domain, modestly; across datasets, clearly. On the five-genre Eerola subset the pipeline
reaches 0.417 against 0.377 for the best direct representation, ahead of all six direct
representations and the PCA-8 control in all ten repeats and, for the routes re-run under
four further master seeds, under every seed. The margin is significant on eight genres under both tests (against AST and VGGish, the
only direct representations run there), but on five genres only under
the film bootstrap ($p=0.042$) and not under the stricter Nadeau--Bengio test ($p=0.151$),
and it disappears when per-genre decision thresholds are tuned on the training data. Across
datasets the picture is different: trained on Eerola and applied to Blockbuster without any
further fitting, the emotion route reaches 0.511 against 0.407 for direct audio ($+0.106$,
$p=0.018$) and 0.421 for the PCA-8 control ($+0.092$, $p=0.032$), against a random-guess
floor of 0.292; the reverse direction agrees ($+0.050$, $p<0.001$), and the gap vanishes
when both datasets are in the training data ($+0.004$, $p=0.850$). The caveats are that the
zero-shot result is a single split resting on one representation, VGGish, and that the
advantage over the PCA-8 control, while positive in every design, is significant only in the
strict zero-shot run ($p=0.032$) and, in-domain, under the film bootstrap alone ($p=0.029$;
Nadeau--Bengio $p=0.122$).

\paragraph{RQ3: Which emotions discriminate each genre most strongly?}
Fear is the master discriminator, with the strongest single genre--emotion association in
the dataset (Horror, $d=+0.75$), and valence is the second axis, separating the light genres
from the threatening ones. Fear defines Horror in both datasets. The ablation study adds the
caveat that ``strongest'' does not mean ``necessary'': fear alone matches all eight emotions
as a classifier input (0.433 against 0.426), and so does the pair of valence and energy,
because the genre signal lies on one low-dimensional axis that almost any pair of emotions
recovers. The discrete emotions are kept because they name what a score is doing, not
because the classifier needs them.

\paragraph{Synthesis.}
Taken together, the results support a specific claim rather than a general one. Routing
film audio through emotion costs nothing in accuracy compared with the strongest direct
representation, and it buys two things: a model whose predictions can be read in terms of
named emotions and checked against film-scoring practice, and a representation that
survives a change of dataset far better than the embedding it is computed from. The modest
and threshold-dependent in-domain margin is part of this finding, not a blemish on it. It is
what one should expect if the value of the emotion intermediate lies in generalisation: when
training and test data come from the same distribution, a rich embedding can learn most of
what emotion provides; when they do not, the emotion axes, which mean the same thing in every
dataset, carry over and the embedding's do not.

\section{Future Work}
\label{sec:futurework}

The following directions are listed in descending order of expected value.

\paragraph{More films.} This is the single remaining lever on the in-domain claims. The
learning curves are still rising, and for every borderline five-genre comparison the
smallest $p$-value the Nadeau--Bengio test could reach with unlimited repetition is above
0.05, so no further computation on the present films can settle them. A larger dataset
could, and it would give the rare genres enough films to be learned at all.

\paragraph{Emotion-annotated film music at scale, with audio on both sides.} The transfer
result rests on VGGish alone, because the second dataset distributes no audio, and on
emotions that are predicted rather than rated on that dataset. A second dataset of film
music with audio and emotion ratings would allow every representation to be carried across,
the predicted emotions to be validated directly, and the Stage-1 regressor to be trained on
more than one listener population. Tuning the regressor itself, which was left at a fixed
configuration while only the genre classifier was tuned, belongs to the same step.

\paragraph{A threshold-free and calibrated evaluation.} The in-domain advantage depends on
the decision threshold, and the zero-shot comparison was run only at the default one. Two
experiments would close this gap: a zero-shot comparison with per-genre thresholds tuned on
the source dataset, and an evaluation of probability calibration alongside the ranking
metrics already reported. Together they would separate what the emotion route contributes to
\emph{ranking} genres from what it contributes to \emph{deciding} them.

\paragraph{Per-cue and temporal modelling.} Every clip is currently summarised by one
time-averaged embedding, and a Blockbuster film by the average of its per-cue emotions.
Applying the emotion model per cue already proved better than averaging the embedding first.
The natural next step is to model how emotion unfolds within a cue and across the cues of a
film -- a horror score is defined as much by its build-up and release as by its average
fear -- and to learn which cues matter most for a film's genre.

\paragraph{Cue-level genre labels and induced emotion.} Film-level genre labels are the
largest source of label noise identified in Chapter~\ref{ch:discussion}; labels assigned
per cue or per scene would remove it. With a suitably annotated dataset, emotion
\emph{induced} in the listener could also be compared with the perceived emotion modelled
here.

\paragraph{Box office, separated from budget.} On the Blockbuster films the emotions of
the score tracked gross only as far as they tracked the production budget. Whether any
property of film music relates to success \emph{beyond} the scale of production could only
be asked with a design that holds the budget fixed -- for example by comparing films of
similar budget and genre -- on a much larger sample, and even then only as an exploratory,
non-causal question.

One direction is deliberately absent from this list: a larger neural network. The direct
representations evaluated here, spanning very different architectures and training data,
are statistically indistinguishable from one another, and the genre signal they need to
recover is low-dimensional. More model capacity is the one change the evidence of this
thesis suggests will not help; more and better-labelled data is the one it suggests will.
